\documentclass[11pt,a4paper]{article}
\usepackage[utf8]{inputenc}
\usepackage[T1]{fontenc}
\usepackage[catalan]{babel}
\usepackage{lmodern}
\usepackage[sfdefault,lining]{FiraSans}
\usepackage[a4paper,top=2.3cm,bottom=2.2cm,left=1.8cm,right=1.8cm,headheight=1.3cm,headsep=0.4cm,footskip=1cm]{geometry}
\usepackage{amsmath,amssymb}
\usepackage{icomma}
\usepackage{graphicx}
\usepackage[table]{xcolor}
\usepackage{tikz}
\usepackage{fontawesome5}
\usepackage{ccicons}
\usepackage{fancyhdr}
\usepackage{lastpage}
\usepackage{tasks}
\usepackage{needspace}
\usepackage[most]{tcolorbox}
\usepackage{array}
\usepackage[hidelinks,pdftitle={Fitxa d'exercicis de la unitat 10: Funcions. 2n d'ESO},pdfsubject={Matemàtiques 2n ESO},pdfkeywords={funcions, variable independent, variable dependent, taules de valors, gràfics, equació d'una funció, funció de proporcionalitat, pendent, funció lineal, ordenada en l'origen, funció constant, interpretació de gràfics}]{hyperref}

\definecolor{navy}{HTML}{1B2748}
\definecolor{accent}{HTML}{E8603A}
\definecolor{teal}{HTML}{1F9E90}
\definecolor{soft}{HTML}{EEF1F8}
\definecolor{muted}{HTML}{5D6685}

\setlength{\parindent}{0pt}
\setlength{\parskip}{3pt}
\settasks{label=\alph*), label-format=\bfseries\color{navy}, label-width=1.3em, item-indent=1.8em, column-sep=0.6em, after-item-skip=5pt, before-skip=3pt, after-skip=4pt}

% ---------- capçalera i peu ----------
\newcommand{\llicencia}{\href{https://creativecommons.org/licenses/by-nc-sa/4.0/deed.ca}{\ccbyncsa}\ \ Llicència Creative Commons Reconeixement-NoComercial-CompartirIgual 4.0 Internacional (CC BY-NC-SA 4.0)}
\pagestyle{fancy}
\fancyhf{}
\fancyhead[L]{\raisebox{-0.35\height}{\includegraphics[height=0.95cm]{logo-gva.png}}}
\fancyhead[C]{\small\color{muted}Fitxa d'exercicis · Unitat 10: Funcions · 2n d'ESO}
\fancyhead[R]{\raisebox{-0.35\height}{\includegraphics[height=0.95cm]{logo-vermellar.png}}}
\fancyfoot[L]{\scriptsize\color{muted}\llicencia}
\fancyfoot[R]{\small\color{muted}\thepage}
\renewcommand{\headrulewidth}{0.4pt}
\renewcommand{\headrule}{\color{muted!40}\hrule height \headrulewidth}
\fancypagestyle{primera}{\fancyhf{}\renewcommand{\headrulewidth}{0pt}
  \fancyfoot[L]{\scriptsize\color{muted}\llicencia}
  \fancyfoot[R]{\small\color{muted}\thepage}}

% ---------- nivells i exercicis ----------
\newcommand{\estrelles}[1]{{\color{accent}\ifcase#1\or\faStar\faStar[regular]\faStar[regular]\or\faStar\faStar\faStar[regular]\or\faStar\faStar\faStar\fi}}
\newcounter{exer}
\newcommand{\ex}[1]{\par\needspace{4\baselineskip}\medskip\stepcounter{exer}\noindent{\color{navy}\bfseries\theexer.}\ \ {\small\estrelles{#1}}\ \ }

\newcommand{\seccio}[2]{\par\needspace{10\baselineskip}\bigskip
  \begin{tcolorbox}[colback=navy,colframe=navy,coltext=white,arc=2mm,boxrule=0pt,left=3mm,right=3mm,top=1.2mm,bottom=1.2mm]
  {\large\bfseries #1\quad #2}
  \end{tcolorbox}\vspace{-2pt}}

\newcolumntype{C}[1]{>{\centering\arraybackslash}p{#1}}

% quadre de recordatori (només a la fitxa d'exercicis)
\newcommand{\recorda}[1]{\par\smallskip\begin{tcolorbox}[colback=accent!7,colframe=accent!55,arc=2mm,boxrule=0.5pt,left=2.5mm,right=2.5mm,top=1mm,bottom=1mm]\small{\color{accent}\faLightbulb[regular]}\ \textbf{Recorda:}\ #1\end{tcolorbox}}
\newcommand{\res}[1]{#1}\newcommand{\rest}[1]{#1}

\def\fmneed{14}
% =====================================================================
% Capçalera FentMates de les fitxes: un color per a cada curs.
% Els colors són els de la pàgina inicial (LEVELS a assets/js/materials.js);
% el logo ve de _eines/logo/pdf/ (python3 _eines/logo/genera.py).
%
% Ús, al final del preàmbul de la unitat:
%   % =====================================================================
% Capçalera FentMates de les fitxes: un color per a cada curs.
% Els colors són els de la pàgina inicial (LEVELS a assets/js/materials.js);
% el logo ve de _eines/logo/pdf/ (python3 _eines/logo/genera.py).
%
% Ús, al final del preàmbul de la unitat:
%   % =====================================================================
% Capçalera FentMates de les fitxes: un color per a cada curs.
% Els colors són els de la pàgina inicial (LEVELS a assets/js/materials.js);
% el logo ve de _eines/logo/pdf/ (python3 _eines/logo/genera.py).
%
% Ús, al final del preàmbul de la unitat:
%   \input{../comu/fentmates.tex}
%   \fmcurs{1eso}
% i, al començament del document:
%   \capcalera{Matemàtiques}{1r ESO}{Fitxa d'exercicis · Unitat 2}{Divisibilitat}
% Opcional abans de \input: \def\fmneed{14} (espai mínim abans de cada secció).
% =====================================================================
\makeatletter
\def\fm@logodir{../../logo/pdf}
\providecommand{\fmneed}{10}
\def\fm@cursnom{}

% clau del curs, color i color del text damunt del color
\newcommand{\fm@curs}[3]{\@namedef{fm@color@#1}{#2}\@namedef{fm@text@#1}{#3}}
\fm@curs{1eso}{E6B84A}{navy}
\fm@curs{2eso}{3B82F6}{white}
\fm@curs{3eso}{FF7A50}{white}
\fm@curs{4eso-a}{A78BFA}{white}
\fm@curs{4eso-b}{6D28D9}{white}
\fm@curs{1batx-mi}{9B1C1C}{white}
\fm@curs{1batx-ccss}{EF6B6B}{white}
\fm@curs{2batx-mii}{1E6B47}{white}
\fm@curs{2batx-ccss}{2BB3A3}{white}

\newcommand{\fmcurs}[1]{%
  \@ifundefined{fm@color@#1}{\PackageError{fentmates}{Curs desconegut: #1}{}}{}%
  \edef\fm@tmp{\noexpand\definecolor{curs}{HTML}{\@nameuse{fm@color@#1}}%
               \noexpand\colorlet{curstext}{\@nameuse{fm@text@#1}}}\fm@tmp
  \colorlet{cursfosc}{curs!72!black}%
  \def\fm@logo{\fm@logodir/vertical-#1.pdf}%
  \def\fm@marca{\fm@logodir/marca-#1.pdf}}

% ---------- capçalera de la primera pàgina ----------
% #1 matèria (petit), #2 curs, #3 tipus de document i unitat, #4 títol
\newcommand{\capcalera}[4]{%
  \gdef\fm@cursnom{#2}%
  \par\noindent
  \begin{tcolorbox}[colback=curs,colframe=curs,boxrule=0pt,arc=2.2mm,width=3.5cm,height=2.15cm,
      valign=center,left=3mm,right=2mm,top=1mm,bottom=1mm,nobeforeafter,box align=center]
    \color{curstext}{\footnotesize #1}\par\vspace{1pt}{\fontsize{19}{21}\selectfont\bfseries #2\par}%
  \end{tcolorbox}\hspace{5mm}%
  \begin{minipage}[c]{\dimexpr\linewidth-3.5cm-5mm-2.9cm\relax}
    {\color{cursfosc}\small\bfseries\MakeUppercase{#3}}\par\vspace{2pt}
    {\color{navy}\fontsize{20}{23}\selectfont\bfseries #4\par}
  \end{minipage}\hfill
  \begin{minipage}[c]{2.6cm}\raggedleft\includegraphics[width=2.6cm]{\fm@logo}\end{minipage}\par
  \vspace{2pt}}

% ---------- barres de secció del color del curs ----------
\renewcommand{\seccio}[2]{\par\needspace{\fmneed\baselineskip}\bigskip
  \begin{tcolorbox}[colback=curs,colframe=curs,coltext=curstext,arc=2mm,boxrule=0pt,left=3mm,right=3mm,top=1.2mm,bottom=1.2mm]
  {\large\bfseries #1\quad #2}
  \end{tcolorbox}\vspace{-2pt}}

% ---------- estrelles de dificultat del color del curs (a Batxillerat no se'n posen) ----------
\@ifundefined{estrelles}{}{%
  \renewcommand{\estrelles}[1]{{\color{curs}\ifcase#1\or\faStar\faStar[regular]\faStar[regular]\or\faStar\faStar\faStar[regular]\or\faStar\faStar\faStar\fi}}}

% ---------- capçalera i peu de les pàgines ----------
% primera pàgina: logos institucionals al peu
\fancypagestyle{primera}{\fancyhf{}\renewcommand{\headrulewidth}{0pt}
  \fancyfoot[L]{\raisebox{-0.32\height}{\includegraphics[height=0.78cm]{logo-gva.png}}\hspace{5mm}%
    \raisebox{-0.32\height}{\includegraphics[height=0.78cm]{logo-vermellar.png}}}
  \fancyfoot[R]{\scriptsize\color{muted}\href{https://creativecommons.org/licenses/by-nc-sa/4.0/deed.ca}{\ccbyncsa}\ \ CC BY-NC-SA 4.0 · fentmates.org\qquad\small\thepage}}
% la resta: el símbol a l'esquerra i el curs a la dreta, amb el títol al centre com abans
\fancyhead[L]{\raisebox{-0.3\height}{\includegraphics[height=0.62cm]{\fm@marca}}}
\fancyhead[R]{\small\bfseries\color{cursfosc}\fm@cursnom}
\renewcommand{\headrule}{\color{curs!45}\hrule height \headrulewidth}
\makeatother

%   \fmcurs{1eso}
% i, al començament del document:
%   \capcalera{Matemàtiques}{1r ESO}{Fitxa d'exercicis · Unitat 2}{Divisibilitat}
% Opcional abans de \input: \def\fmneed{14} (espai mínim abans de cada secció).
% =====================================================================
\makeatletter
\def\fm@logodir{../../logo/pdf}
\providecommand{\fmneed}{10}
\def\fm@cursnom{}

% clau del curs, color i color del text damunt del color
\newcommand{\fm@curs}[3]{\@namedef{fm@color@#1}{#2}\@namedef{fm@text@#1}{#3}}
\fm@curs{1eso}{E6B84A}{navy}
\fm@curs{2eso}{3B82F6}{white}
\fm@curs{3eso}{FF7A50}{white}
\fm@curs{4eso-a}{A78BFA}{white}
\fm@curs{4eso-b}{6D28D9}{white}
\fm@curs{1batx-mi}{9B1C1C}{white}
\fm@curs{1batx-ccss}{EF6B6B}{white}
\fm@curs{2batx-mii}{1E6B47}{white}
\fm@curs{2batx-ccss}{2BB3A3}{white}

\newcommand{\fmcurs}[1]{%
  \@ifundefined{fm@color@#1}{\PackageError{fentmates}{Curs desconegut: #1}{}}{}%
  \edef\fm@tmp{\noexpand\definecolor{curs}{HTML}{\@nameuse{fm@color@#1}}%
               \noexpand\colorlet{curstext}{\@nameuse{fm@text@#1}}}\fm@tmp
  \colorlet{cursfosc}{curs!72!black}%
  \def\fm@logo{\fm@logodir/vertical-#1.pdf}%
  \def\fm@marca{\fm@logodir/marca-#1.pdf}}

% ---------- capçalera de la primera pàgina ----------
% #1 matèria (petit), #2 curs, #3 tipus de document i unitat, #4 títol
\newcommand{\capcalera}[4]{%
  \gdef\fm@cursnom{#2}%
  \par\noindent
  \begin{tcolorbox}[colback=curs,colframe=curs,boxrule=0pt,arc=2.2mm,width=3.5cm,height=2.15cm,
      valign=center,left=3mm,right=2mm,top=1mm,bottom=1mm,nobeforeafter,box align=center]
    \color{curstext}{\footnotesize #1}\par\vspace{1pt}{\fontsize{19}{21}\selectfont\bfseries #2\par}%
  \end{tcolorbox}\hspace{5mm}%
  \begin{minipage}[c]{\dimexpr\linewidth-3.5cm-5mm-2.9cm\relax}
    {\color{cursfosc}\small\bfseries\MakeUppercase{#3}}\par\vspace{2pt}
    {\color{navy}\fontsize{20}{23}\selectfont\bfseries #4\par}
  \end{minipage}\hfill
  \begin{minipage}[c]{2.6cm}\raggedleft\includegraphics[width=2.6cm]{\fm@logo}\end{minipage}\par
  \vspace{2pt}}

% ---------- barres de secció del color del curs ----------
\renewcommand{\seccio}[2]{\par\needspace{\fmneed\baselineskip}\bigskip
  \begin{tcolorbox}[colback=curs,colframe=curs,coltext=curstext,arc=2mm,boxrule=0pt,left=3mm,right=3mm,top=1.2mm,bottom=1.2mm]
  {\large\bfseries #1\quad #2}
  \end{tcolorbox}\vspace{-2pt}}

% ---------- estrelles de dificultat del color del curs (a Batxillerat no se'n posen) ----------
\@ifundefined{estrelles}{}{%
  \renewcommand{\estrelles}[1]{{\color{curs}\ifcase#1\or\faStar\faStar[regular]\faStar[regular]\or\faStar\faStar\faStar[regular]\or\faStar\faStar\faStar\fi}}}

% ---------- capçalera i peu de les pàgines ----------
% primera pàgina: logos institucionals al peu
\fancypagestyle{primera}{\fancyhf{}\renewcommand{\headrulewidth}{0pt}
  \fancyfoot[L]{\raisebox{-0.32\height}{\includegraphics[height=0.78cm]{logo-gva.png}}\hspace{5mm}%
    \raisebox{-0.32\height}{\includegraphics[height=0.78cm]{logo-vermellar.png}}}
  \fancyfoot[R]{\scriptsize\color{muted}\href{https://creativecommons.org/licenses/by-nc-sa/4.0/deed.ca}{\ccbyncsa}\ \ CC BY-NC-SA 4.0 · fentmates.org\qquad\small\thepage}}
% la resta: el símbol a l'esquerra i el curs a la dreta, amb el títol al centre com abans
\fancyhead[L]{\raisebox{-0.3\height}{\includegraphics[height=0.62cm]{\fm@marca}}}
\fancyhead[R]{\small\bfseries\color{cursfosc}\fm@cursnom}
\renewcommand{\headrule}{\color{curs!45}\hrule height \headrulewidth}
\makeatother

%   \fmcurs{1eso}
% i, al començament del document:
%   \capcalera{Matemàtiques}{1r ESO}{Fitxa d'exercicis · Unitat 2}{Divisibilitat}
% Opcional abans de \input: \def\fmneed{14} (espai mínim abans de cada secció).
% =====================================================================
\makeatletter
\def\fm@logodir{../../logo/pdf}
\providecommand{\fmneed}{10}
\def\fm@cursnom{}

% clau del curs, color i color del text damunt del color
\newcommand{\fm@curs}[3]{\@namedef{fm@color@#1}{#2}\@namedef{fm@text@#1}{#3}}
\fm@curs{1eso}{E6B84A}{navy}
\fm@curs{2eso}{3B82F6}{white}
\fm@curs{3eso}{FF7A50}{white}
\fm@curs{4eso-a}{A78BFA}{white}
\fm@curs{4eso-b}{6D28D9}{white}
\fm@curs{1batx-mi}{9B1C1C}{white}
\fm@curs{1batx-ccss}{EF6B6B}{white}
\fm@curs{2batx-mii}{1E6B47}{white}
\fm@curs{2batx-ccss}{2BB3A3}{white}

\newcommand{\fmcurs}[1]{%
  \@ifundefined{fm@color@#1}{\PackageError{fentmates}{Curs desconegut: #1}{}}{}%
  \edef\fm@tmp{\noexpand\definecolor{curs}{HTML}{\@nameuse{fm@color@#1}}%
               \noexpand\colorlet{curstext}{\@nameuse{fm@text@#1}}}\fm@tmp
  \colorlet{cursfosc}{curs!72!black}%
  \def\fm@logo{\fm@logodir/vertical-#1.pdf}%
  \def\fm@marca{\fm@logodir/marca-#1.pdf}}

% ---------- capçalera de la primera pàgina ----------
% #1 matèria (petit), #2 curs, #3 tipus de document i unitat, #4 títol
\newcommand{\capcalera}[4]{%
  \gdef\fm@cursnom{#2}%
  \par\noindent
  \begin{tcolorbox}[colback=curs,colframe=curs,boxrule=0pt,arc=2.2mm,width=3.5cm,height=2.15cm,
      valign=center,left=3mm,right=2mm,top=1mm,bottom=1mm,nobeforeafter,box align=center]
    \color{curstext}{\footnotesize #1}\par\vspace{1pt}{\fontsize{19}{21}\selectfont\bfseries #2\par}%
  \end{tcolorbox}\hspace{5mm}%
  \begin{minipage}[c]{\dimexpr\linewidth-3.5cm-5mm-2.9cm\relax}
    {\color{cursfosc}\small\bfseries\MakeUppercase{#3}}\par\vspace{2pt}
    {\color{navy}\fontsize{20}{23}\selectfont\bfseries #4\par}
  \end{minipage}\hfill
  \begin{minipage}[c]{2.6cm}\raggedleft\includegraphics[width=2.6cm]{\fm@logo}\end{minipage}\par
  \vspace{2pt}}

% ---------- barres de secció del color del curs ----------
\renewcommand{\seccio}[2]{\par\needspace{\fmneed\baselineskip}\bigskip
  \begin{tcolorbox}[colback=curs,colframe=curs,coltext=curstext,arc=2mm,boxrule=0pt,left=3mm,right=3mm,top=1.2mm,bottom=1.2mm]
  {\large\bfseries #1\quad #2}
  \end{tcolorbox}\vspace{-2pt}}

% ---------- estrelles de dificultat del color del curs (a Batxillerat no se'n posen) ----------
\@ifundefined{estrelles}{}{%
  \renewcommand{\estrelles}[1]{{\color{curs}\ifcase#1\or\faStar\faStar[regular]\faStar[regular]\or\faStar\faStar\faStar[regular]\or\faStar\faStar\faStar\fi}}}

% ---------- capçalera i peu de les pàgines ----------
% primera pàgina: logos institucionals al peu
\fancypagestyle{primera}{\fancyhf{}\renewcommand{\headrulewidth}{0pt}
  \fancyfoot[L]{\raisebox{-0.32\height}{\includegraphics[height=0.78cm]{logo-gva.png}}\hspace{5mm}%
    \raisebox{-0.32\height}{\includegraphics[height=0.78cm]{logo-vermellar.png}}}
  \fancyfoot[R]{\scriptsize\color{muted}\href{https://creativecommons.org/licenses/by-nc-sa/4.0/deed.ca}{\ccbyncsa}\ \ CC BY-NC-SA 4.0 · fentmates.org\qquad\small\thepage}}
% la resta: el símbol a l'esquerra i el curs a la dreta, amb el títol al centre com abans
\fancyhead[L]{\raisebox{-0.3\height}{\includegraphics[height=0.62cm]{\fm@marca}}}
\fancyhead[R]{\small\bfseries\color{cursfosc}\fm@cursnom}
\renewcommand{\headrule}{\color{curs!45}\hrule height \headrulewidth}
\makeatother

\fmcurs{2eso}

\begin{document}
\thispagestyle{primera}

\capcalera{Matemàtiques}{2n ESO}{Fitxa d'exercicis · Unitat 10}{Funcions}
\vspace{10pt}
Nom i cognoms: \rule{8.4cm}{0.4pt}\hfill Grup: \rule{1.6cm}{0.4pt}\hfill Data: \rule{2.2cm}{0.4pt}
\vspace{8pt}

\begin{tcolorbox}[colback=soft,colframe=soft,arc=2mm,left=3mm,right=3mm,top=1.5mm,bottom=1.5mm]
\small Exercicis graduats per nivells de dificultat:\quad \estrelles{1}\ bàsic \qquad \estrelles{2}\ mitjà \qquad \estrelles{3}\ avançat.\\[2pt] Resol-los al quadern mostrant tots els passos. Per a dibuixar gràfics fes servir paper quadriculat i escriu el nom i la unitat de cada eix. Les seccions segueixen l'ordre dels apunts de la unitat (10.1--10.8).
\end{tcolorbox}

% =====================================================
\seccio{1}{Concepte de funció}

\ex{1} En cada situació, digues quina és la variable independent ($x$) i quina la dependent ($y$):\begin{tasks}(2)
  \task El preu d'una pizza i el nombre d'ingredients.
  \task L'àrea d'un cercle i el seu radi.
  \task Els minuts parlats per telèfon i el preu.
  \task La temperatura d'un forn i el temps que fa que està encés.
\end{tasks}\recorda{La variable \textbf{independent} $x$ és la que triem o la que passa sola; la \textbf{dependent} $y$ és la que depén de l'altra. Una relació és una \textbf{funció} si a cada valor de $x$ li correspon \textbf{un únic} valor de $y$.}

\ex{1} Digues si aquestes relacions són funcions i explica per què:\begin{tasks}(1)
  \task A cada alumne de la classe, el seu dia de naixement.
  \task A cada nombre natural, els seus múltiples.
  \task A cada quantitat de litres de gasolina, el preu que paguem.
  \task A cada edat, les persones de la classe que la tenen.
\end{tasks}

\ex{1} Quines d'aquestes taules són d'una funció? Raona-ho.\par\smallskip\begin{center}$\begin{array}{|c|cccc|}\hline x&1&2&3&4\\\hline y&3&6&9&12\\\hline\end{array}$\hspace{0.7cm}$\begin{array}{|c|cccc|}\hline x&1&1&2&3\\\hline y&2&5&4&7\\\hline\end{array}$\hspace{0.7cm}$\begin{array}{|c|cccc|}\hline x&-2&-1&1&2\\\hline y&4&1&1&4\\\hline\end{array}$\end{center}

\ex{2} Quins d'aquests gràfics són d'una funció? Fes la prova de la recta vertical.\par\smallskip\begin{center}\begin{tabular}{c}\begin{tikzpicture}[x=0.3cm,y=0.3cm,baseline=(current bounding box.center),line cap=round,line join=round]
\draw[muted!30,line width=0.3pt] (-4,-3)--(-4,3);
\draw[muted!30,line width=0.3pt] (-3,-3)--(-3,3);
\draw[muted!30,line width=0.3pt] (-2,-3)--(-2,3);
\draw[muted!30,line width=0.3pt] (-1,-3)--(-1,3);
\draw[muted!30,line width=0.3pt] (0,-3)--(0,3);
\draw[muted!30,line width=0.3pt] (1,-3)--(1,3);
\draw[muted!30,line width=0.3pt] (2,-3)--(2,3);
\draw[muted!30,line width=0.3pt] (3,-3)--(3,3);
\draw[muted!30,line width=0.3pt] (4,-3)--(4,3);
\draw[muted!30,line width=0.3pt] (-4,-3)--(4,-3);
\draw[muted!30,line width=0.3pt] (-4,-2)--(4,-2);
\draw[muted!30,line width=0.3pt] (-4,-1)--(4,-1);
\draw[muted!30,line width=0.3pt] (-4,0)--(4,0);
\draw[muted!30,line width=0.3pt] (-4,1)--(4,1);
\draw[muted!30,line width=0.3pt] (-4,2)--(4,2);
\draw[muted!30,line width=0.3pt] (-4,3)--(4,3);
\draw[-latex,muted,line width=0.6pt] (-4,0)--(4.3,0) node[right,font=\scriptsize,text=navy]{x};
\draw[-latex,muted,line width=0.6pt] (0,-3)--(0,3.3) node[above,font=\scriptsize,text=navy]{y};
\node[font=\tiny,text=muted,below,inner sep=1.5pt] at (-2,0) {$-2$};
\node[font=\tiny,text=muted,below,inner sep=1.5pt] at (2,0) {$2$};
\node[font=\tiny,text=muted,left,inner sep=1.5pt] at (0,-2) {$-2$};
\node[font=\tiny,text=muted,left,inner sep=1.5pt] at (0,2) {$2$};
\node[font=\tiny,text=muted,below left,inner sep=1.5pt] at (0,0) {$0$};
\draw[accent,line width=1.1pt] (-3.500,-2.000)--(-1.500,1.500)--(0.500,0.000)--(3.500,2.500);
\end{tikzpicture}\\[1pt]\small A\end{tabular}\hspace{0.7cm}\begin{tabular}{c}\begin{tikzpicture}[x=0.3cm,y=0.3cm,baseline=(current bounding box.center),line cap=round,line join=round]
\draw[muted!30,line width=0.3pt] (-4,-3)--(-4,3);
\draw[muted!30,line width=0.3pt] (-3,-3)--(-3,3);
\draw[muted!30,line width=0.3pt] (-2,-3)--(-2,3);
\draw[muted!30,line width=0.3pt] (-1,-3)--(-1,3);
\draw[muted!30,line width=0.3pt] (0,-3)--(0,3);
\draw[muted!30,line width=0.3pt] (1,-3)--(1,3);
\draw[muted!30,line width=0.3pt] (2,-3)--(2,3);
\draw[muted!30,line width=0.3pt] (3,-3)--(3,3);
\draw[muted!30,line width=0.3pt] (4,-3)--(4,3);
\draw[muted!30,line width=0.3pt] (-4,-3)--(4,-3);
\draw[muted!30,line width=0.3pt] (-4,-2)--(4,-2);
\draw[muted!30,line width=0.3pt] (-4,-1)--(4,-1);
\draw[muted!30,line width=0.3pt] (-4,0)--(4,0);
\draw[muted!30,line width=0.3pt] (-4,1)--(4,1);
\draw[muted!30,line width=0.3pt] (-4,2)--(4,2);
\draw[muted!30,line width=0.3pt] (-4,3)--(4,3);
\draw[-latex,muted,line width=0.6pt] (-4,0)--(4.3,0) node[right,font=\scriptsize,text=navy]{x};
\draw[-latex,muted,line width=0.6pt] (0,-3)--(0,3.3) node[above,font=\scriptsize,text=navy]{y};
\node[font=\tiny,text=muted,below,inner sep=1.5pt] at (-2,0) {$-2$};
\node[font=\tiny,text=muted,below,inner sep=1.5pt] at (2,0) {$2$};
\node[font=\tiny,text=muted,left,inner sep=1.5pt] at (0,-2) {$-2$};
\node[font=\tiny,text=muted,left,inner sep=1.5pt] at (0,2) {$2$};
\node[font=\tiny,text=muted,below left,inner sep=1.5pt] at (0,0) {$0$};
\draw[accent,line width=1.1pt] (2.200,0.000)--(2.167,0.382)--(2.067,0.752)--(1.905,1.100)--(1.685,1.414)--(1.414,1.685)--(1.100,1.905)--(0.752,2.067)--(0.382,2.167)--(0.000,2.200)--(-0.382,2.167)--(-0.752,2.067)--(-1.100,1.905)--(-1.414,1.685)--(-1.685,1.414)--(-1.905,1.100)--(-2.067,0.752)--(-2.167,0.382)--(-2.200,0.000)--(-2.167,-0.382)--(-2.067,-0.752)--(-1.905,-1.100)--(-1.685,-1.414)--(-1.414,-1.685)--(-1.100,-1.905)--(-0.752,-2.067)--(-0.382,-2.167)--(-0.000,-2.200)--(0.382,-2.167)--(0.752,-2.067)--(1.100,-1.905)--(1.414,-1.685)--(1.685,-1.414)--(1.905,-1.100)--(2.067,-0.752)--(2.167,-0.382)--(2.200,-0.000);
\end{tikzpicture}\\[1pt]\small B\end{tabular}\hspace{0.7cm}\begin{tabular}{c}\begin{tikzpicture}[x=0.3cm,y=0.3cm,baseline=(current bounding box.center),line cap=round,line join=round]
\draw[muted!30,line width=0.3pt] (-4,-3)--(-4,3);
\draw[muted!30,line width=0.3pt] (-3,-3)--(-3,3);
\draw[muted!30,line width=0.3pt] (-2,-3)--(-2,3);
\draw[muted!30,line width=0.3pt] (-1,-3)--(-1,3);
\draw[muted!30,line width=0.3pt] (0,-3)--(0,3);
\draw[muted!30,line width=0.3pt] (1,-3)--(1,3);
\draw[muted!30,line width=0.3pt] (2,-3)--(2,3);
\draw[muted!30,line width=0.3pt] (3,-3)--(3,3);
\draw[muted!30,line width=0.3pt] (4,-3)--(4,3);
\draw[muted!30,line width=0.3pt] (-4,-3)--(4,-3);
\draw[muted!30,line width=0.3pt] (-4,-2)--(4,-2);
\draw[muted!30,line width=0.3pt] (-4,-1)--(4,-1);
\draw[muted!30,line width=0.3pt] (-4,0)--(4,0);
\draw[muted!30,line width=0.3pt] (-4,1)--(4,1);
\draw[muted!30,line width=0.3pt] (-4,2)--(4,2);
\draw[muted!30,line width=0.3pt] (-4,3)--(4,3);
\draw[-latex,muted,line width=0.6pt] (-4,0)--(4.3,0) node[right,font=\scriptsize,text=navy]{x};
\draw[-latex,muted,line width=0.6pt] (0,-3)--(0,3.3) node[above,font=\scriptsize,text=navy]{y};
\node[font=\tiny,text=muted,below,inner sep=1.5pt] at (-2,0) {$-2$};
\node[font=\tiny,text=muted,below,inner sep=1.5pt] at (2,0) {$2$};
\node[font=\tiny,text=muted,left,inner sep=1.5pt] at (0,-2) {$-2$};
\node[font=\tiny,text=muted,left,inner sep=1.5pt] at (0,2) {$2$};
\node[font=\tiny,text=muted,below left,inner sep=1.5pt] at (0,0) {$0$};
\draw[accent,line width=1.1pt] (3.750,-2.500)--(3.260,-2.400)--(2.790,-2.300)--(2.340,-2.200)--(1.910,-2.100)--(1.500,-2.000)--(1.110,-1.900)--(0.740,-1.800)--(0.390,-1.700)--(0.060,-1.600)--(-0.250,-1.500)--(-0.540,-1.400)--(-0.810,-1.300)--(-1.060,-1.200)--(-1.290,-1.100)--(-1.500,-1.000)--(-1.690,-0.900)--(-1.860,-0.800)--(-2.010,-0.700)--(-2.140,-0.600)--(-2.250,-0.500)--(-2.340,-0.400)--(-2.410,-0.300)--(-2.460,-0.200)--(-2.490,-0.100)--(-2.500,0.000)--(-2.490,0.100)--(-2.460,0.200)--(-2.410,0.300)--(-2.340,0.400)--(-2.250,0.500)--(-2.140,0.600)--(-2.010,0.700)--(-1.860,0.800)--(-1.690,0.900)--(-1.500,1.000)--(-1.290,1.100)--(-1.060,1.200)--(-0.810,1.300)--(-0.540,1.400)--(-0.250,1.500)--(0.060,1.600)--(0.390,1.700)--(0.740,1.800)--(1.110,1.900)--(1.500,2.000)--(1.910,2.100)--(2.340,2.200)--(2.790,2.300)--(3.260,2.400)--(3.750,2.500);
\end{tikzpicture}\\[1pt]\small C\end{tabular}\hspace{0.7cm}\begin{tabular}{c}\begin{tikzpicture}[x=0.3cm,y=0.3cm,baseline=(current bounding box.center),line cap=round,line join=round]
\draw[muted!30,line width=0.3pt] (-4,-3)--(-4,3);
\draw[muted!30,line width=0.3pt] (-3,-3)--(-3,3);
\draw[muted!30,line width=0.3pt] (-2,-3)--(-2,3);
\draw[muted!30,line width=0.3pt] (-1,-3)--(-1,3);
\draw[muted!30,line width=0.3pt] (0,-3)--(0,3);
\draw[muted!30,line width=0.3pt] (1,-3)--(1,3);
\draw[muted!30,line width=0.3pt] (2,-3)--(2,3);
\draw[muted!30,line width=0.3pt] (3,-3)--(3,3);
\draw[muted!30,line width=0.3pt] (4,-3)--(4,3);
\draw[muted!30,line width=0.3pt] (-4,-3)--(4,-3);
\draw[muted!30,line width=0.3pt] (-4,-2)--(4,-2);
\draw[muted!30,line width=0.3pt] (-4,-1)--(4,-1);
\draw[muted!30,line width=0.3pt] (-4,0)--(4,0);
\draw[muted!30,line width=0.3pt] (-4,1)--(4,1);
\draw[muted!30,line width=0.3pt] (-4,2)--(4,2);
\draw[muted!30,line width=0.3pt] (-4,3)--(4,3);
\draw[-latex,muted,line width=0.6pt] (-4,0)--(4.3,0) node[right,font=\scriptsize,text=navy]{x};
\draw[-latex,muted,line width=0.6pt] (0,-3)--(0,3.3) node[above,font=\scriptsize,text=navy]{y};
\node[font=\tiny,text=muted,below,inner sep=1.5pt] at (-2,0) {$-2$};
\node[font=\tiny,text=muted,below,inner sep=1.5pt] at (2,0) {$2$};
\node[font=\tiny,text=muted,left,inner sep=1.5pt] at (0,-2) {$-2$};
\node[font=\tiny,text=muted,left,inner sep=1.5pt] at (0,2) {$2$};
\node[font=\tiny,text=muted,below left,inner sep=1.5pt] at (0,0) {$0$};
\draw[accent,line width=1.1pt] (-4.000,-1.500)--(4.000,-1.500);
\end{tikzpicture}\\[1pt]\small D\end{tabular}\end{center}\recorda{Si alguna recta vertical talla el gràfic en dos punts o més, \textbf{no} és una funció.}

\ex{2} Una impressora imprimeix 20 pàgines cada minut. Dona aquesta funció amb una taula (de 0 a 4 minuts), amb una fórmula i amb un gràfic. Unim els punts?

\ex{3} Els rectangles d'àrea 24\,cm$^2$ poden tindre bases diferents. a) Completa la taula de l'altura $h$ segons la base $b$. b) És una funció? c) Quina fórmula té? d) Quan la base es duplica, què passa amb l'altura?\par\smallskip\begin{center}$\begin{array}{|c|C{0.75cm}C{0.75cm}C{0.75cm}C{0.75cm}C{0.75cm}C{0.75cm}C{0.75cm}C{0.75cm}|}\hline \text{b}&1&2&3&4&6&8&12&24\\\hline \text{h}&\phantom{00}&\phantom{00}&\phantom{00}&\phantom{00}&\phantom{00}&\phantom{00}&\phantom{00}&\phantom{00}\\\hline\end{array}$\end{center}


% =====================================================
\seccio{2}{Funcions donades per taules de valors}

\ex{1} Les mandarines valen 1,20\,€ el quilo. Completa la taula del preu ($y$) segons el pes ($x$) i digues si en el gràfic unim els punts.\par\smallskip\begin{center}$\begin{array}{|c|C{0.75cm}C{0.75cm}C{0.75cm}C{0.75cm}C{0.75cm}|}\hline \text{kg}&0,5&1&2&3&4\\\hline \text{€}&\phantom{00}&\phantom{00}&\phantom{00}&\phantom{00}&\phantom{00}\\\hline\end{array}$\end{center}\recorda{Cada columna de la taula és un punt $(x,\ y)$ del gràfic. Unim els punts si també tenen sentit els valors intermedis; si només hi ha valors solts (persones, entrades…), no els unim.}

\ex{1} Digues si en el gràfic d'aquestes funcions s'han d'unir els punts:\begin{tasks}(1)
  \task Nombre de germans d'un alumne → nombre de regals que fa per Nadal.
  \task Temps que fa que corre un atleta → distància recorreguda.
  \task Nombre de llibres comprats → preu total.
  \task Litres d'aigua d'un dipòsit → altura de l'aigua.
\end{tasks}

\ex{1} La taula dona la temperatura d'un dia cada 3 hores. Representa-la i respon: a) A quina hora va fer més calor? I més fred? b) Entre quines hores va pujar la temperatura? c) Quina temperatura aproximada feia a les 10 h?\par\smallskip\begin{center}$\begin{array}{|c|ccccccccc|}\hline \text{h}&0&3&6&9&12&15&18&21&24\\\hline \text{°C}&8&6&7&13&18&20&16&11&9\\\hline\end{array}$\end{center}

\ex{2} Troba la fórmula de la funció que dona cada taula:\par\smallskip\begin{center}$\begin{array}{|c|ccccc|}\hline x&0&1&2&3&4\\\hline y&1&4&7&10&13\\\hline\end{array}$\hspace{0.7cm}$\begin{array}{|c|ccccc|}\hline x&0&1&2&3&4\\\hline y&20&17&14&11&8\\\hline\end{array}$\end{center}\recorda{Si $x$ va d'1 en 1 i $y$ augmenta (o disminueix) sempre igual, la fórmula és $y=mx+n$: $m$ és el que canvia $y$ cada vegada i $n$ és el valor de $y$ per a $x=0$.}

\ex{2} Completa la taula de la funció $y=2x-5$:\par\smallskip\begin{center}$\begin{array}{|c|cccccc|}\hline x&-1&0&2&\phantom{00}&\phantom{00}&10\\\hline y&\phantom{00}&\phantom{00}&\phantom{00}&3&0&\phantom{00}\\\hline\end{array}$\end{center}

\ex{2} Una piscina es buida. La taula dona l'altura de l'aigua (cm) segons les hores. a) Quanta aigua baixa cada hora? b) Quina altura hi havia a les 5 h? c) Quan quedaran 20\,cm? d) Escriu la fórmula.\par\smallskip\begin{center}$\begin{array}{|c|ccccc|}\hline \text{h}&0&2&4&6&8\\\hline \text{cm}&90&70&50&30&10\\\hline\end{array}$\end{center}

\ex{3} Fes la taula de l'àrea $A$ d'un quadrat segons el costat $l$ per a $l=1,\ 2,\ 3,\ 4,\ 5$ i representa-la. És una recta? Quina àrea té un quadrat de 2,5\,cm de costat?

\ex{3} Amb escuradents fem una filera de triangles:\par\smallskip\begin{center}\begin{tikzpicture}[x=0.5cm,y=0.5cm,line width=1pt,accent]\draw (0,0)--(0.5,0.87)--(1,0);\draw (0.5,0.87)--(1.5,0.87);\draw (1,0)--(1.5,0.87)--(2,0);\draw (1.5,0.87)--(2.5,0.87);\draw (2,0)--(2.5,0.87)--(3,0);\draw (2.5,0.87)--(3.5,0.87);\draw (3,0)--(3.5,0.87)--(4,0);\end{tikzpicture}\end{center}a) Completa la taula del nombre d'escuradents $y$ segons el nombre de triangles $x$. b) Troba'n la fórmula. c) Quants escuradents calen per a 20 triangles? d) Quants triangles fem amb 61 escuradents?\par\smallskip\begin{center}$\begin{array}{|c|C{0.75cm}C{0.75cm}C{0.75cm}C{0.75cm}|}\hline \text{triangles}&1&2&3&4\\\hline \text{escuradents}&\phantom{00}&\phantom{00}&\phantom{00}&\phantom{00}\\\hline\end{array}$\end{center}


% =====================================================
\seccio{3}{Funcions donades per la seua equació}

\ex{1} Calcula el valor de $y=4x-3$ per a:\begin{tasks}(4)
  \task $x=-2$
  \task $x=0$
  \task $x=\frac12$
  \task $x=3$
\end{tasks}\recorda{Substituïm la $x$ pel seu valor (entre parèntesis si és negatiu) i fem les operacions. El punt $(x,\ y)$ obtingut és del gràfic.}

\ex{1} Fes una taula de valors de $y=-x+2$ per a $x$ de $-2$ a $4$ i representa-la.

\ex{1} Fes la taula de valors de $y=x^2-2$ per a $x=-3,\ -2,\ \ldots,\ 3$ i dibuixa'n el gràfic.

\ex{2} Quins d'aquests punts són del gràfic de $y=3x-1$?\begin{tasks}(4)
  \task $A(2,\ 5)$
  \task $B(-1,\ -4)$
  \task $C(0,\ 1)$
  \task $D\left(\frac13,\ 0\right)$
\end{tasks}

\ex{2} Troba el valor de $x$ en cada cas:\begin{tasks}(1)
  \task Per a quin valor de $x$ la funció $y=5x+2$ val 17?
  \task Per a quin valor de $x$ la funció $y=-2x+3$ val 9?
  \task Per a quins valors de $x$ la funció $y=x^2$ val 16?
\end{tasks}

\ex{2} Fes la taula de $y=\frac{12}{x}$ per a $x=1,\ 2,\ 3,\ 4,\ 6,\ 12$ i dibuixa'n el gràfic. Quant val $y$ per a $x=0$?

\ex{3} Un rectangle té 16\,m de perímetre i base $x$. a) Comprova que l'altura és $8-x$ i escriu l'equació de l'àrea $A$. b) Fes la taula per a $x=1,\ \ldots,\ 7$ i el gràfic. c) Quines dimensions donen l'àrea màxima?

\ex{3} La distància que necessita un cotxe per a frenar (en m) es pot calcular amb $d=\frac{v^2}{100}$, on $v$ és la velocitat en km/h. a) Fes la taula per a 30, 50, 80, 100 i 120\,km/h. b) Si la velocitat es duplica de 50 a 100\,km/h, la distància també es duplica? c) Quina distància necessita a 90\,km/h?


% =====================================================
\seccio{4}{Funció de proporcionalitat: $y=mx$}

\ex{1} Quines taules corresponen a una funció de proporcionalitat? Escriu-ne l'equació.\par\smallskip\begin{center}$\begin{array}{|c|ccc|}\hline x&1&2&4\\\hline y&5&10&20\\\hline\end{array}$\hspace{0.7cm}$\begin{array}{|c|ccc|}\hline x&2&3&5\\\hline y&6&8&12\\\hline\end{array}$\hspace{0.7cm}$\begin{array}{|c|ccc|}\hline x&2&4&6\\\hline y&-1&-2&-3\\\hline\end{array}$\end{center}\recorda{En una funció de proporcionalitat $y=mx$, el quocient $\frac yx$ és sempre $m$. El gràfic és una recta que passa per l'origen.}

\ex{1} Representa en els mateixos eixos $y=2x$, $y=-3x$ i $y=\frac x3$. Quines són creixents? Quina és la més inclinada?

\ex{1} Escriu l'equació de la recta que passa per l'origen i pel punt:\begin{tasks}(3)
  \task $(3,\ 12)$
  \task $(-2,\ 5)$
  \task $(4,\ -2)$
\end{tasks}

\ex{2} Una tela costa 4,50\,€ el metre. Escriu la funció del preu $y$ segons els metres $x$, fes una taula i el gràfic. Quants metres comprem amb 27\,€?

\ex{2} Un cotxe gasta 6 litres de gasolina cada 100\,km. a) Escriu la funció dels litres $y$ que gasta segons els km $x$. b) Quant gasta en 350\,km? c) Quants km pot fer amb 45 litres?

\ex{2} Sense dibuixar-les, digues si aquestes rectes són creixents o decreixents i ordena-les de menys a més inclinada: $y=-5x$, \ $y=\frac12x$, \ $y=2x$, \ $y=-x$.

\ex{3} Un dia, 1 euro equival a 0,85 lliures. a) Escriu la funció que passa d'euros ($x$) a lliures ($y$). b) Quantes lliures són 60\,€? c) Quants euros són 34 lliures? d) Escriu també la funció que passa de lliures a euros. És de proporcionalitat?


% =====================================================
\seccio{5}{Pendent d'una recta}

\ex{1} Calcula el pendent de la recta que passa per cada parell de punts i digues si és creixent, decreixent o constant:\begin{tasks}(2)
  \task $(1,\ 2)$ i $(4,\ 11)$
  \task $(-2,\ 5)$ i $(1,\ -1)$
  \task $(0,\ 3)$ i $(6,\ 7)$
  \task $(-3,\ 4)$ i $(5,\ 4)$
\end{tasks}\recorda{$m=\dfrac{y_2-y_1}{x_2-x_1}$: és el que augmenta $y$ quan $x$ augmenta 1. Si $m>0$ la recta puja; si $m<0$, baixa; si $m=0$, és horitzontal.}

\ex{1} Llig el pendent de cada recta en el gràfic (busca dos punts que caiguen en la quadrícula):\par\smallskip\begin{center}\begin{tabular}{c}\begin{tikzpicture}[x=0.3cm,y=0.3cm,baseline=(current bounding box.center),line cap=round,line join=round]
\draw[muted!30,line width=0.3pt] (-4,-4)--(-4,4);
\draw[muted!30,line width=0.3pt] (-3,-4)--(-3,4);
\draw[muted!30,line width=0.3pt] (-2,-4)--(-2,4);
\draw[muted!30,line width=0.3pt] (-1,-4)--(-1,4);
\draw[muted!30,line width=0.3pt] (0,-4)--(0,4);
\draw[muted!30,line width=0.3pt] (1,-4)--(1,4);
\draw[muted!30,line width=0.3pt] (2,-4)--(2,4);
\draw[muted!30,line width=0.3pt] (3,-4)--(3,4);
\draw[muted!30,line width=0.3pt] (4,-4)--(4,4);
\draw[muted!30,line width=0.3pt] (-4,-4)--(4,-4);
\draw[muted!30,line width=0.3pt] (-4,-3)--(4,-3);
\draw[muted!30,line width=0.3pt] (-4,-2)--(4,-2);
\draw[muted!30,line width=0.3pt] (-4,-1)--(4,-1);
\draw[muted!30,line width=0.3pt] (-4,0)--(4,0);
\draw[muted!30,line width=0.3pt] (-4,1)--(4,1);
\draw[muted!30,line width=0.3pt] (-4,2)--(4,2);
\draw[muted!30,line width=0.3pt] (-4,3)--(4,3);
\draw[muted!30,line width=0.3pt] (-4,4)--(4,4);
\draw[-latex,muted,line width=0.6pt] (-4,0)--(4.3,0) node[right,font=\scriptsize,text=navy]{x};
\draw[-latex,muted,line width=0.6pt] (0,-4)--(0,4.3) node[above,font=\scriptsize,text=navy]{y};
\node[font=\tiny,text=muted,below,inner sep=1.5pt] at (-3,0) {$-3$};
\node[font=\tiny,text=muted,below,inner sep=1.5pt] at (-2,0) {$-2$};
\node[font=\tiny,text=muted,below,inner sep=1.5pt] at (-1,0) {$-1$};
\node[font=\tiny,text=muted,below,inner sep=1.5pt] at (1,0) {$1$};
\node[font=\tiny,text=muted,below,inner sep=1.5pt] at (2,0) {$2$};
\node[font=\tiny,text=muted,below,inner sep=1.5pt] at (3,0) {$3$};
\node[font=\tiny,text=muted,left,inner sep=1.5pt] at (0,-3) {$-3$};
\node[font=\tiny,text=muted,left,inner sep=1.5pt] at (0,-2) {$-2$};
\node[font=\tiny,text=muted,left,inner sep=1.5pt] at (0,-1) {$-1$};
\node[font=\tiny,text=muted,left,inner sep=1.5pt] at (0,1) {$1$};
\node[font=\tiny,text=muted,left,inner sep=1.5pt] at (0,2) {$2$};
\node[font=\tiny,text=muted,left,inner sep=1.5pt] at (0,3) {$3$};
\node[font=\tiny,text=muted,below left,inner sep=1.5pt] at (0,0) {$0$};
\draw[accent,line width=1.1pt] (-1.500,-4.000)--(2.500,4.000);
\end{tikzpicture}\\[1pt]\small a)\end{tabular}\hspace{0.7cm}\begin{tabular}{c}\begin{tikzpicture}[x=0.3cm,y=0.3cm,baseline=(current bounding box.center),line cap=round,line join=round]
\draw[muted!30,line width=0.3pt] (-4,-4)--(-4,4);
\draw[muted!30,line width=0.3pt] (-3,-4)--(-3,4);
\draw[muted!30,line width=0.3pt] (-2,-4)--(-2,4);
\draw[muted!30,line width=0.3pt] (-1,-4)--(-1,4);
\draw[muted!30,line width=0.3pt] (0,-4)--(0,4);
\draw[muted!30,line width=0.3pt] (1,-4)--(1,4);
\draw[muted!30,line width=0.3pt] (2,-4)--(2,4);
\draw[muted!30,line width=0.3pt] (3,-4)--(3,4);
\draw[muted!30,line width=0.3pt] (4,-4)--(4,4);
\draw[muted!30,line width=0.3pt] (-4,-4)--(4,-4);
\draw[muted!30,line width=0.3pt] (-4,-3)--(4,-3);
\draw[muted!30,line width=0.3pt] (-4,-2)--(4,-2);
\draw[muted!30,line width=0.3pt] (-4,-1)--(4,-1);
\draw[muted!30,line width=0.3pt] (-4,0)--(4,0);
\draw[muted!30,line width=0.3pt] (-4,1)--(4,1);
\draw[muted!30,line width=0.3pt] (-4,2)--(4,2);
\draw[muted!30,line width=0.3pt] (-4,3)--(4,3);
\draw[muted!30,line width=0.3pt] (-4,4)--(4,4);
\draw[-latex,muted,line width=0.6pt] (-4,0)--(4.3,0) node[right,font=\scriptsize,text=navy]{x};
\draw[-latex,muted,line width=0.6pt] (0,-4)--(0,4.3) node[above,font=\scriptsize,text=navy]{y};
\node[font=\tiny,text=muted,below,inner sep=1.5pt] at (-3,0) {$-3$};
\node[font=\tiny,text=muted,below,inner sep=1.5pt] at (-2,0) {$-2$};
\node[font=\tiny,text=muted,below,inner sep=1.5pt] at (-1,0) {$-1$};
\node[font=\tiny,text=muted,below,inner sep=1.5pt] at (1,0) {$1$};
\node[font=\tiny,text=muted,below,inner sep=1.5pt] at (2,0) {$2$};
\node[font=\tiny,text=muted,below,inner sep=1.5pt] at (3,0) {$3$};
\node[font=\tiny,text=muted,left,inner sep=1.5pt] at (0,-3) {$-3$};
\node[font=\tiny,text=muted,left,inner sep=1.5pt] at (0,-2) {$-2$};
\node[font=\tiny,text=muted,left,inner sep=1.5pt] at (0,-1) {$-1$};
\node[font=\tiny,text=muted,left,inner sep=1.5pt] at (0,1) {$1$};
\node[font=\tiny,text=muted,left,inner sep=1.5pt] at (0,2) {$2$};
\node[font=\tiny,text=muted,left,inner sep=1.5pt] at (0,3) {$3$};
\node[font=\tiny,text=muted,below left,inner sep=1.5pt] at (0,0) {$0$};
\draw[accent,line width=1.1pt] (-4.000,4.000)--(4.000,0.000);
\end{tikzpicture}\\[1pt]\small b)\end{tabular}\hspace{0.7cm}\begin{tabular}{c}\begin{tikzpicture}[x=0.3cm,y=0.3cm,baseline=(current bounding box.center),line cap=round,line join=round]
\draw[muted!30,line width=0.3pt] (-4,-4)--(-4,4);
\draw[muted!30,line width=0.3pt] (-3,-4)--(-3,4);
\draw[muted!30,line width=0.3pt] (-2,-4)--(-2,4);
\draw[muted!30,line width=0.3pt] (-1,-4)--(-1,4);
\draw[muted!30,line width=0.3pt] (0,-4)--(0,4);
\draw[muted!30,line width=0.3pt] (1,-4)--(1,4);
\draw[muted!30,line width=0.3pt] (2,-4)--(2,4);
\draw[muted!30,line width=0.3pt] (3,-4)--(3,4);
\draw[muted!30,line width=0.3pt] (4,-4)--(4,4);
\draw[muted!30,line width=0.3pt] (-4,-4)--(4,-4);
\draw[muted!30,line width=0.3pt] (-4,-3)--(4,-3);
\draw[muted!30,line width=0.3pt] (-4,-2)--(4,-2);
\draw[muted!30,line width=0.3pt] (-4,-1)--(4,-1);
\draw[muted!30,line width=0.3pt] (-4,0)--(4,0);
\draw[muted!30,line width=0.3pt] (-4,1)--(4,1);
\draw[muted!30,line width=0.3pt] (-4,2)--(4,2);
\draw[muted!30,line width=0.3pt] (-4,3)--(4,3);
\draw[muted!30,line width=0.3pt] (-4,4)--(4,4);
\draw[-latex,muted,line width=0.6pt] (-4,0)--(4.3,0) node[right,font=\scriptsize,text=navy]{x};
\draw[-latex,muted,line width=0.6pt] (0,-4)--(0,4.3) node[above,font=\scriptsize,text=navy]{y};
\node[font=\tiny,text=muted,below,inner sep=1.5pt] at (-3,0) {$-3$};
\node[font=\tiny,text=muted,below,inner sep=1.5pt] at (-2,0) {$-2$};
\node[font=\tiny,text=muted,below,inner sep=1.5pt] at (-1,0) {$-1$};
\node[font=\tiny,text=muted,below,inner sep=1.5pt] at (1,0) {$1$};
\node[font=\tiny,text=muted,below,inner sep=1.5pt] at (2,0) {$2$};
\node[font=\tiny,text=muted,below,inner sep=1.5pt] at (3,0) {$3$};
\node[font=\tiny,text=muted,left,inner sep=1.5pt] at (0,-3) {$-3$};
\node[font=\tiny,text=muted,left,inner sep=1.5pt] at (0,-2) {$-2$};
\node[font=\tiny,text=muted,left,inner sep=1.5pt] at (0,-1) {$-1$};
\node[font=\tiny,text=muted,left,inner sep=1.5pt] at (0,1) {$1$};
\node[font=\tiny,text=muted,left,inner sep=1.5pt] at (0,2) {$2$};
\node[font=\tiny,text=muted,left,inner sep=1.5pt] at (0,3) {$3$};
\node[font=\tiny,text=muted,below left,inner sep=1.5pt] at (0,0) {$0$};
\draw[accent,line width=1.1pt] (-4.000,-4.000)--(4.000,2.000);
\end{tikzpicture}\\[1pt]\small c)\end{tabular}\hspace{0.7cm}\begin{tabular}{c}\begin{tikzpicture}[x=0.3cm,y=0.3cm,baseline=(current bounding box.center),line cap=round,line join=round]
\draw[muted!30,line width=0.3pt] (-4,-4)--(-4,4);
\draw[muted!30,line width=0.3pt] (-3,-4)--(-3,4);
\draw[muted!30,line width=0.3pt] (-2,-4)--(-2,4);
\draw[muted!30,line width=0.3pt] (-1,-4)--(-1,4);
\draw[muted!30,line width=0.3pt] (0,-4)--(0,4);
\draw[muted!30,line width=0.3pt] (1,-4)--(1,4);
\draw[muted!30,line width=0.3pt] (2,-4)--(2,4);
\draw[muted!30,line width=0.3pt] (3,-4)--(3,4);
\draw[muted!30,line width=0.3pt] (4,-4)--(4,4);
\draw[muted!30,line width=0.3pt] (-4,-4)--(4,-4);
\draw[muted!30,line width=0.3pt] (-4,-3)--(4,-3);
\draw[muted!30,line width=0.3pt] (-4,-2)--(4,-2);
\draw[muted!30,line width=0.3pt] (-4,-1)--(4,-1);
\draw[muted!30,line width=0.3pt] (-4,0)--(4,0);
\draw[muted!30,line width=0.3pt] (-4,1)--(4,1);
\draw[muted!30,line width=0.3pt] (-4,2)--(4,2);
\draw[muted!30,line width=0.3pt] (-4,3)--(4,3);
\draw[muted!30,line width=0.3pt] (-4,4)--(4,4);
\draw[-latex,muted,line width=0.6pt] (-4,0)--(4.3,0) node[right,font=\scriptsize,text=navy]{x};
\draw[-latex,muted,line width=0.6pt] (0,-4)--(0,4.3) node[above,font=\scriptsize,text=navy]{y};
\node[font=\tiny,text=muted,below,inner sep=1.5pt] at (-3,0) {$-3$};
\node[font=\tiny,text=muted,below,inner sep=1.5pt] at (-2,0) {$-2$};
\node[font=\tiny,text=muted,below,inner sep=1.5pt] at (-1,0) {$-1$};
\node[font=\tiny,text=muted,below,inner sep=1.5pt] at (1,0) {$1$};
\node[font=\tiny,text=muted,below,inner sep=1.5pt] at (2,0) {$2$};
\node[font=\tiny,text=muted,below,inner sep=1.5pt] at (3,0) {$3$};
\node[font=\tiny,text=muted,left,inner sep=1.5pt] at (0,-3) {$-3$};
\node[font=\tiny,text=muted,left,inner sep=1.5pt] at (0,-2) {$-2$};
\node[font=\tiny,text=muted,left,inner sep=1.5pt] at (0,-1) {$-1$};
\node[font=\tiny,text=muted,left,inner sep=1.5pt] at (0,1) {$1$};
\node[font=\tiny,text=muted,left,inner sep=1.5pt] at (0,2) {$2$};
\node[font=\tiny,text=muted,left,inner sep=1.5pt] at (0,3) {$3$};
\node[font=\tiny,text=muted,below left,inner sep=1.5pt] at (0,0) {$0$};
\draw[accent,line width=1.1pt] (-0.333,4.000)--(2.333,-4.000);
\end{tikzpicture}\\[1pt]\small d)\end{tabular}\end{center}

\ex{2} Dibuixa la recta que passa per $A(1,\ 2)$ i té pendent 3, i la recta que passa per $B(0,\ 4)$ i té pendent $-\frac12$.

\ex{2} Una carretera té un pendent del 12\,\%: puja 12\,m per cada 100\,m en horitzontal. a) Quant puja en 3\,km? b) Una rampa té un pendent del 8\,\% i ha de pujar 0,6\,m. Quina longitud horitzontal ha de tindre?

\ex{2} Comprova amb els pendents si els punts estan alineats:\begin{tasks}(1)
  \task $A(0,\ 1)$, $B(2,\ 7)$, $C(5,\ 16)$
  \task $P(-1,\ 3)$, $Q(1,\ 2)$, $R(5,\ -1)$
\end{tasks}

\ex{2} Sense dibuixar-les, ordena les rectes de menys a més inclinada: $r$: passa per $(0,\ 0)$ i $(2,\ 6)$; \ $s$: passa per $(1,\ 1)$ i $(5,\ 3)$; \ $t$: passa per $(0,\ 4)$ i $(1,\ 0)$.

\ex{3} a) Quin valor ha de tindre $k$ perquè els punts $(0,\ 1)$, $(2,\ 5)$ i $(4,\ k)$ estiguen alineats? b) Quin valor ha de tindre $a$ perquè el pendent de la recta que passa per $(a,\ 3)$ i $(2,\ 7)$ siga 2?


% =====================================================
\seccio{6}{Funcions lineals: $y=mx+n$}

\ex{1} Escriu el pendent $m$ i l'ordenada en l'origen $n$ de cada recta:\begin{tasks}(5)
  \task $y=3x-2$
  \task $y=-x+5$
  \task $y=4-2x$
  \task $y=\frac x2+3$
  \task $y=-6x$
\end{tasks}\recorda{En $y=mx+n$, $m$ és el \textbf{pendent} i $n$ l'\textbf{ordenada en l'origen}: la recta talla l'eix $Y$ en $(0,\ n)$. Les rectes amb el mateix pendent són paral·leles.}

\ex{1} Representa en els mateixos eixos $y=x+1$ i $y=-x+3$. En quin punt es tallen?

\ex{1} Escriu l'equació de les rectes $r$, $s$ i $t$:\par\smallskip\begin{center}\begin{tikzpicture}[x=0.42cm,y=0.42cm,baseline=(current bounding box.center),line cap=round,line join=round]
\draw[muted!30,line width=0.3pt] (-4,-4)--(-4,4);
\draw[muted!30,line width=0.3pt] (-3,-4)--(-3,4);
\draw[muted!30,line width=0.3pt] (-2,-4)--(-2,4);
\draw[muted!30,line width=0.3pt] (-1,-4)--(-1,4);
\draw[muted!30,line width=0.3pt] (0,-4)--(0,4);
\draw[muted!30,line width=0.3pt] (1,-4)--(1,4);
\draw[muted!30,line width=0.3pt] (2,-4)--(2,4);
\draw[muted!30,line width=0.3pt] (3,-4)--(3,4);
\draw[muted!30,line width=0.3pt] (4,-4)--(4,4);
\draw[muted!30,line width=0.3pt] (-4,-4)--(4,-4);
\draw[muted!30,line width=0.3pt] (-4,-3)--(4,-3);
\draw[muted!30,line width=0.3pt] (-4,-2)--(4,-2);
\draw[muted!30,line width=0.3pt] (-4,-1)--(4,-1);
\draw[muted!30,line width=0.3pt] (-4,0)--(4,0);
\draw[muted!30,line width=0.3pt] (-4,1)--(4,1);
\draw[muted!30,line width=0.3pt] (-4,2)--(4,2);
\draw[muted!30,line width=0.3pt] (-4,3)--(4,3);
\draw[muted!30,line width=0.3pt] (-4,4)--(4,4);
\draw[-latex,muted,line width=0.6pt] (-4,0)--(4.3,0) node[right,font=\scriptsize,text=navy]{x};
\draw[-latex,muted,line width=0.6pt] (0,-4)--(0,4.3) node[above,font=\scriptsize,text=navy]{y};
\node[font=\tiny,text=muted,below,inner sep=1.5pt] at (-3,0) {$-3$};
\node[font=\tiny,text=muted,below,inner sep=1.5pt] at (-2,0) {$-2$};
\node[font=\tiny,text=muted,below,inner sep=1.5pt] at (-1,0) {$-1$};
\node[font=\tiny,text=muted,below,inner sep=1.5pt] at (1,0) {$1$};
\node[font=\tiny,text=muted,below,inner sep=1.5pt] at (2,0) {$2$};
\node[font=\tiny,text=muted,below,inner sep=1.5pt] at (3,0) {$3$};
\node[font=\tiny,text=muted,left,inner sep=1.5pt] at (0,-3) {$-3$};
\node[font=\tiny,text=muted,left,inner sep=1.5pt] at (0,-2) {$-2$};
\node[font=\tiny,text=muted,left,inner sep=1.5pt] at (0,-1) {$-1$};
\node[font=\tiny,text=muted,left,inner sep=1.5pt] at (0,1) {$1$};
\node[font=\tiny,text=muted,left,inner sep=1.5pt] at (0,2) {$2$};
\node[font=\tiny,text=muted,left,inner sep=1.5pt] at (0,3) {$3$};
\node[font=\tiny,text=muted,below left,inner sep=1.5pt] at (0,0) {$0$};
\draw[accent,line width=1.1pt] (-1.500,-4.000)--(2.500,4.000);
\draw[teal,line width=1.1pt] (-3.000,4.000)--(4.000,-0.667);
\draw[navy,line width=1.1pt] (-4.000,-3.000)--(4.000,-3.000);
\node[font=\scriptsize,right] at (2.2,3.6) {$r$};
\node[font=\scriptsize,below] at (-3.6,3.6) {$s$};
\node[font=\scriptsize,above] at (3.6,-3) {$t$};
\end{tikzpicture}\end{center}

\ex{2} Escriu l'equació de la recta que:\begin{tasks}(1)
  \task té pendent 3 i passa pel punt $(1,\ -2)$;
  \task té pendent $-\frac12$ i passa pel punt $(4,\ 1)$;
  \task és paral·lela a $y=-2x+1$ i passa per $(3,\ 0)$.
\end{tasks}

\ex{2} Troba l'equació de la recta que passa per cada parell de punts:\begin{tasks}(2)
  \task $(-1,\ 4)$ i $(2,\ -2)$
  \task $(0,\ -3)$ i $(4,\ 3)$
\end{tasks}

\ex{2} Comprova si els punts $A(4,\ 7)$, $B(-1,\ -8)$ i $C(0{,}5;\ -3{,}5)$ són de la recta $y=3x-5$.

\ex{3} La recta $y=2x-6$ talla els eixos en dos punts. Troba'ls i calcula l'àrea del triangle que forma la recta amb els eixos.

\ex{3} L'aigua es congela a 0\,°C, que són 32\,°F (graus Fahrenheit), i bull a 100\,°C, que són 212\,°F. La relació és lineal. a) Troba la fórmula $F=mC+n$. b) Quants °F són 25\,°C? c) Quants °C són 50\,°F?


% =====================================================
\seccio{7}{Funcions constants: $y=k$}

\ex{1} Representa en els mateixos eixos les funcions $y=2$, $y=-3$ i $y=0$. Quin pendent tenen?\recorda{Una \textbf{funció constant} $y=k$ val sempre el mateix: el gràfic és una recta \textbf{horitzontal} i el pendent és 0. Una recta vertical $x=a$ \textbf{no} és una funció.}

\ex{1} Escriu l'equació de la recta horitzontal que passa per $(4,\ -1)$ i de la recta vertical que passa per $(-2,\ 5)$. Quina és una funció?

\ex{1} Classifica les funcions en de proporcionalitat, lineals (amb $n\neq0$), constants o cap d'aquestes:\begin{tasks}(4)
  \task $y=-4x$
  \task $y=7$
  \task $y=2x+1$
  \task $y=x^2+1$
  \task $y=\frac{x}{5}$
  \task $y=3-x$
  \task $y=\frac{6}{x}$
  \task $y=-2$
\end{tasks}

\ex{2} Enviar un paquet costa 4\,€ fins a 2\,kg, 7\,€ de 2 a 5\,kg i 10\,€ de 5 a 10\,kg. a) Dibuixa el gràfic. b) Quant costa enviar un paquet de 3,5\,kg? c) Què ix més barat: enviar dos paquets de 3\,kg o un de 6\,kg?

\ex{2} Una companyia telefònica ofereix una tarifa plana de 30\,€ al mes i una altra cobra 0,08\,€ per minut. a) Escriu les dues funcions. b) Quant paguem amb cada una si parlem 200 minuts? c) A partir de quants minuts convé la tarifa plana?

\ex{3} La taula correspon a una funció lineal. Troba'n l'equació i explica com és el gràfic. Quin valor tindrà $y$ per a $x=100$?\par\smallskip\begin{center}$\begin{array}{|c|ccccc|}\hline x&-2&-1&0&1&2\\\hline y&5&5&5&5&5\\\hline\end{array}$\end{center}


% =====================================================
\seccio{8}{Interpretació de gràfics i problemes}

\ex{1} El gràfic mostra la distància a casa d'un excursionista. a) A quina distància estava a la 1\,h i mitja? b) Quant de temps va descansar? c) Quina distància màxima es va allunyar? d) Quants km va recórrer en total? e) Va més de pressa a l'anada o a la tornada?\par\smallskip\begin{center}\begin{tikzpicture}[x=0.9cm,y=0.3cm,baseline=(current bounding box.center),line cap=round,line join=round]
\draw[muted!30,line width=0.3pt] (0,0)--(0,11.5);
\draw[muted!30,line width=0.3pt] (1,0)--(1,11.5);
\draw[muted!30,line width=0.3pt] (2,0)--(2,11.5);
\draw[muted!30,line width=0.3pt] (3,0)--(3,11.5);
\draw[muted!30,line width=0.3pt] (4,0)--(4,11.5);
\draw[muted!30,line width=0.3pt] (5,0)--(5,11.5);
\draw[muted!30,line width=0.3pt] (0,0)--(5.6,0);
\draw[muted!30,line width=0.3pt] (0,1)--(5.6,1);
\draw[muted!30,line width=0.3pt] (0,2)--(5.6,2);
\draw[muted!30,line width=0.3pt] (0,3)--(5.6,3);
\draw[muted!30,line width=0.3pt] (0,4)--(5.6,4);
\draw[muted!30,line width=0.3pt] (0,5)--(5.6,5);
\draw[muted!30,line width=0.3pt] (0,6)--(5.6,6);
\draw[muted!30,line width=0.3pt] (0,7)--(5.6,7);
\draw[muted!30,line width=0.3pt] (0,8)--(5.6,8);
\draw[muted!30,line width=0.3pt] (0,9)--(5.6,9);
\draw[muted!30,line width=0.3pt] (0,10)--(5.6,10);
\draw[muted!30,line width=0.3pt] (0,11)--(5.6,11);
\draw[-latex,muted,line width=0.6pt] (0,0)--(5.9,0) node[right,font=\scriptsize,text=navy]{t (h)};
\draw[-latex,muted,line width=0.6pt] (0,0)--(0,11.8) node[above,font=\scriptsize,text=navy]{d (km)};
\node[font=\tiny,text=muted,below,inner sep=1.5pt] at (1,0) {$1$};
\node[font=\tiny,text=muted,below,inner sep=1.5pt] at (2,0) {$2$};
\node[font=\tiny,text=muted,below,inner sep=1.5pt] at (3,0) {$3$};
\node[font=\tiny,text=muted,below,inner sep=1.5pt] at (4,0) {$4$};
\node[font=\tiny,text=muted,below,inner sep=1.5pt] at (5,0) {$5$};
\node[font=\tiny,text=muted,left,inner sep=1.5pt] at (0,2) {$2$};
\node[font=\tiny,text=muted,left,inner sep=1.5pt] at (0,4) {$4$};
\node[font=\tiny,text=muted,left,inner sep=1.5pt] at (0,6) {$6$};
\node[font=\tiny,text=muted,left,inner sep=1.5pt] at (0,8) {$8$};
\node[font=\tiny,text=muted,left,inner sep=1.5pt] at (0,10) {$10$};
\node[font=\tiny,text=muted,below left,inner sep=1.5pt] at (0,0) {$0$};
\draw[accent,line width=1.1pt] (0.000,0.000)--(1.000,6.000)--(2.000,6.000)--(3.000,10.000)--(5.000,0.000);
\fill[navy] (0,0) circle (2.2pt);
\fill[navy] (1,6) circle (2.2pt);
\fill[navy] (2,6) circle (2.2pt);
\fill[navy] (3,10) circle (2.2pt);
\fill[navy] (5,0) circle (2.2pt);
\end{tikzpicture}\end{center}\recorda{Llegim el gràfic d'esquerra a dreta: on puja s'allunya, on és horitzontal està parat i on baixa torna. Com més inclinat, més de pressa.}

\ex{1} El gràfic dona la temperatura d'un dia. Digues on és creixent i on és decreixent, i quines són la temperatura màxima i la mínima i quan es van assolir.\par\smallskip\begin{center}\begin{tikzpicture}[x=0.3cm,y=0.2cm,baseline=(current bounding box.center),line cap=round,line join=round]
\draw[muted!30,line width=0.3pt] (0,0)--(0,20);
\draw[muted!30,line width=0.3pt] (2,0)--(2,20);
\draw[muted!30,line width=0.3pt] (4,0)--(4,20);
\draw[muted!30,line width=0.3pt] (6,0)--(6,20);
\draw[muted!30,line width=0.3pt] (8,0)--(8,20);
\draw[muted!30,line width=0.3pt] (10,0)--(10,20);
\draw[muted!30,line width=0.3pt] (12,0)--(12,20);
\draw[muted!30,line width=0.3pt] (14,0)--(14,20);
\draw[muted!30,line width=0.3pt] (16,0)--(16,20);
\draw[muted!30,line width=0.3pt] (18,0)--(18,20);
\draw[muted!30,line width=0.3pt] (20,0)--(20,20);
\draw[muted!30,line width=0.3pt] (22,0)--(22,20);
\draw[muted!30,line width=0.3pt] (24,0)--(24,20);
\draw[muted!30,line width=0.3pt] (0,0)--(25,0);
\draw[muted!30,line width=0.3pt] (0,2)--(25,2);
\draw[muted!30,line width=0.3pt] (0,4)--(25,4);
\draw[muted!30,line width=0.3pt] (0,6)--(25,6);
\draw[muted!30,line width=0.3pt] (0,8)--(25,8);
\draw[muted!30,line width=0.3pt] (0,10)--(25,10);
\draw[muted!30,line width=0.3pt] (0,12)--(25,12);
\draw[muted!30,line width=0.3pt] (0,14)--(25,14);
\draw[muted!30,line width=0.3pt] (0,16)--(25,16);
\draw[muted!30,line width=0.3pt] (0,18)--(25,18);
\draw[muted!30,line width=0.3pt] (0,20)--(25,20);
\draw[-latex,muted,line width=0.6pt] (0,0)--(25.6,0) node[right,font=\scriptsize,text=navy]{h};
\draw[-latex,muted,line width=0.6pt] (0,0)--(0,20.6) node[above,font=\scriptsize,text=navy]{°C};
\node[font=\tiny,text=muted,below,inner sep=1.5pt] at (4,0) {$4$};
\node[font=\tiny,text=muted,below,inner sep=1.5pt] at (8,0) {$8$};
\node[font=\tiny,text=muted,below,inner sep=1.5pt] at (12,0) {$12$};
\node[font=\tiny,text=muted,below,inner sep=1.5pt] at (16,0) {$16$};
\node[font=\tiny,text=muted,below,inner sep=1.5pt] at (20,0) {$20$};
\node[font=\tiny,text=muted,below,inner sep=1.5pt] at (24,0) {$24$};
\node[font=\tiny,text=muted,left,inner sep=1.5pt] at (0,4) {$4$};
\node[font=\tiny,text=muted,left,inner sep=1.5pt] at (0,8) {$8$};
\node[font=\tiny,text=muted,left,inner sep=1.5pt] at (0,12) {$12$};
\node[font=\tiny,text=muted,left,inner sep=1.5pt] at (0,16) {$16$};
\node[font=\tiny,text=muted,below left,inner sep=1.5pt] at (0,0) {$0$};
\draw[accent,line width=1.1pt] (0.000,10.000)--(4.000,6.000)--(8.000,8.000)--(12.000,16.000)--(16.000,18.000)--(20.000,12.000)--(24.000,9.000);
\fill[navy] (0,10) circle (2.2pt);
\fill[navy] (4,6) circle (2.2pt);
\fill[navy] (8,8) circle (2.2pt);
\fill[navy] (12,16) circle (2.2pt);
\fill[navy] (16,18) circle (2.2pt);
\fill[navy] (20,12) circle (2.2pt);
\fill[navy] (24,9) circle (2.2pt);
\end{tikzpicture}\end{center}

\ex{2} Relaciona cada història amb el seu gràfic (temps, distància a casa):\begin{tasks}(1)
  \task Vaig en bici a casa d'una amiga, m'hi quede una estona i torne.
  \task Isc corrent, pare a descansar i després continue caminant a poc a poc.
  \task Camine sempre a la mateixa velocitat, sense parar.
\end{tasks}\par\smallskip\begin{center}\begin{tabular}{c}\begin{tikzpicture}[x=0.28cm,y=0.28cm,baseline=(current bounding box.center),line cap=round,line join=round]
\draw[muted!30,line width=0.3pt] (0,0)--(0,6.5);
\draw[muted!30,line width=0.3pt] (2,0)--(2,6.5);
\draw[muted!30,line width=0.3pt] (4,0)--(4,6.5);
\draw[muted!30,line width=0.3pt] (6,0)--(6,6.5);
\draw[muted!30,line width=0.3pt] (0,0)--(6.5,0);
\draw[muted!30,line width=0.3pt] (0,2)--(6.5,2);
\draw[muted!30,line width=0.3pt] (0,4)--(6.5,4);
\draw[muted!30,line width=0.3pt] (0,6)--(6.5,6);
\draw[-latex,muted,line width=0.6pt] (0,0)--(7.1,0) node[right,font=\scriptsize,text=navy]{t};
\draw[-latex,muted,line width=0.6pt] (0,0)--(0,7.1) node[above,font=\scriptsize,text=navy]{d};
\node[font=\tiny,text=muted,below,inner sep=1.5pt] at (2,0) {$2$};
\node[font=\tiny,text=muted,below,inner sep=1.5pt] at (4,0) {$4$};
\node[font=\tiny,text=muted,below,inner sep=1.5pt] at (6,0) {$6$};
\node[font=\tiny,text=muted,left,inner sep=1.5pt] at (0,2) {$2$};
\node[font=\tiny,text=muted,left,inner sep=1.5pt] at (0,4) {$4$};
\node[font=\tiny,text=muted,left,inner sep=1.5pt] at (0,6) {$6$};
\node[font=\tiny,text=muted,below left,inner sep=1.5pt] at (0,0) {$0$};
\draw[accent,line width=1.1pt] (0.000,0.000)--(2.000,5.000)--(4.000,5.000)--(6.000,0.000);
\end{tikzpicture}\\[1pt]\small A\end{tabular}\hspace{0.7cm}\begin{tabular}{c}\begin{tikzpicture}[x=0.28cm,y=0.28cm,baseline=(current bounding box.center),line cap=round,line join=round]
\draw[muted!30,line width=0.3pt] (0,0)--(0,6.5);
\draw[muted!30,line width=0.3pt] (2,0)--(2,6.5);
\draw[muted!30,line width=0.3pt] (4,0)--(4,6.5);
\draw[muted!30,line width=0.3pt] (6,0)--(6,6.5);
\draw[muted!30,line width=0.3pt] (0,0)--(6.5,0);
\draw[muted!30,line width=0.3pt] (0,2)--(6.5,2);
\draw[muted!30,line width=0.3pt] (0,4)--(6.5,4);
\draw[muted!30,line width=0.3pt] (0,6)--(6.5,6);
\draw[-latex,muted,line width=0.6pt] (0,0)--(7.1,0) node[right,font=\scriptsize,text=navy]{t};
\draw[-latex,muted,line width=0.6pt] (0,0)--(0,7.1) node[above,font=\scriptsize,text=navy]{d};
\node[font=\tiny,text=muted,below,inner sep=1.5pt] at (2,0) {$2$};
\node[font=\tiny,text=muted,below,inner sep=1.5pt] at (4,0) {$4$};
\node[font=\tiny,text=muted,below,inner sep=1.5pt] at (6,0) {$6$};
\node[font=\tiny,text=muted,left,inner sep=1.5pt] at (0,2) {$2$};
\node[font=\tiny,text=muted,left,inner sep=1.5pt] at (0,4) {$4$};
\node[font=\tiny,text=muted,left,inner sep=1.5pt] at (0,6) {$6$};
\node[font=\tiny,text=muted,below left,inner sep=1.5pt] at (0,0) {$0$};
\draw[accent,line width=1.1pt] (0.000,0.000)--(3.000,3.000)--(6.000,6.000);
\end{tikzpicture}\\[1pt]\small B\end{tabular}\hspace{0.7cm}\begin{tabular}{c}\begin{tikzpicture}[x=0.28cm,y=0.28cm,baseline=(current bounding box.center),line cap=round,line join=round]
\draw[muted!30,line width=0.3pt] (0,0)--(0,6.5);
\draw[muted!30,line width=0.3pt] (2,0)--(2,6.5);
\draw[muted!30,line width=0.3pt] (4,0)--(4,6.5);
\draw[muted!30,line width=0.3pt] (6,0)--(6,6.5);
\draw[muted!30,line width=0.3pt] (0,0)--(6.5,0);
\draw[muted!30,line width=0.3pt] (0,2)--(6.5,2);
\draw[muted!30,line width=0.3pt] (0,4)--(6.5,4);
\draw[muted!30,line width=0.3pt] (0,6)--(6.5,6);
\draw[-latex,muted,line width=0.6pt] (0,0)--(7.1,0) node[right,font=\scriptsize,text=navy]{t};
\draw[-latex,muted,line width=0.6pt] (0,0)--(0,7.1) node[above,font=\scriptsize,text=navy]{d};
\node[font=\tiny,text=muted,below,inner sep=1.5pt] at (2,0) {$2$};
\node[font=\tiny,text=muted,below,inner sep=1.5pt] at (4,0) {$4$};
\node[font=\tiny,text=muted,below,inner sep=1.5pt] at (6,0) {$6$};
\node[font=\tiny,text=muted,left,inner sep=1.5pt] at (0,2) {$2$};
\node[font=\tiny,text=muted,left,inner sep=1.5pt] at (0,4) {$4$};
\node[font=\tiny,text=muted,left,inner sep=1.5pt] at (0,6) {$6$};
\node[font=\tiny,text=muted,below left,inner sep=1.5pt] at (0,0) {$0$};
\draw[accent,line width=1.1pt] (0.000,0.000)--(1.000,5.000)--(3.000,5.000)--(6.000,6.000);
\end{tikzpicture}\\[1pt]\small C\end{tabular}\end{center}

\ex{2} Un taxi cobra 3,50\,€ per la baixada de bandera i 1,10\,€ per km. a) Escriu la funció del preu. b) Quant costa un viatge de 12\,km? c) Quants km hem fet si hem pagat 25,50\,€?\recorda{Moltes situacions tenen una part fixa $n$ i una part que augmenta (o disminueix) sempre igual per cada unitat, $m$: són funcions lineals $y=mx+n$.}

\ex{2} Un dipòsit té 800 litres i se'n buiden 25 cada minut. a) Escriu la funció dels litres que queden. b) Quants litres queden als 12 minuts? c) Quan es buidarà?

\ex{2} La Clara ix de casa a les 10\,h i camina 3\,km en 45 minuts fins a la piscina. Hi està 1\,h 30\,min i torna amb autobús en 15 minuts. Dibuixa el gràfic (hora, distància a casa).

\ex{3} El gimnàs A cobra 30\,€ de matrícula i 15\,€ al mes, i el gimnàs B, 45\,€ de matrícula i 10\,€ al mes. a) Escriu les dues funcions. b) Quant costa cada un si hi anem 6 mesos? c) Per a quants mesos costen el mateix?

\ex{3} Un ciclista ix d'un poble a 18\,km/h. Una hora després n'ix un altre pel mateix camí a 24\,km/h. a) Escriu la distància recorreguda per cada un en funció del temps $t$ que fa que ha eixit el primer. b) Quan i on l'atrapa?

\ex{3} Encenem alhora dues espelmes: una fa 30\,cm i es crema 2\,cm cada hora; l'altra fa 24\,cm i es crema 1,5\,cm cada hora. a) Escriu les dues funcions. b) Quan tindran la mateixa altura? Quina? c) Quina s'acaba abans?

\vfill
\begin{center}\color{muted}\small\faCheckCircle[regular]\ \ Ànim! Comprova els teus resultats amb les eines interactives dels apunts de la unitat.\end{center}

\end{document}
